% LTeX: enabled=false
% language
\usepackage[T1]{fontenc}
\usepackage[utf8]{inputenc}
\usepackage[main=naustrian,english,latin]{babel}


% math commands
\usepackage{amsmath,amssymb}

% bibliography
\usepackage{IEEEtrantools}
\bstctlcite{IEEEexample:BSTcontrol}
\bibliographystyle{IEEEtran}

% insert figures
\usepackage{graphicx}
\pdfminorversion=7

% tables
\usepackage{array}
\renewcommand{\arraystretch}{1.3}

% Listings (for code)
\usepackage{listings}
\usepackage{scrhack}
\lstset{numbers=left,numberstyle=\tiny,stepnumber=5,numbersep=5pt}

% Headers and footers
\usepackage{scrlayer-scrpage}
\automark{chapter}
\cfoot*{\footnotesize\pagemark}
\chead{}
\ohead{\footnotesize\sc\headmark}


% Hyperlinks
\usepackage[colorlinks=true, allcolors=black, bookmarksnumbered=true]{hyperref}
\usepackage{bookmark}

% Helvetica font
\usepackage{helvet}
\renewcommand{\familydefault}{phv}

% spacing
\usepackage{setspace}
\onehalfspacing
\setlength{\parindent}{0em}
\setlength{\parskip}{1.5ex plus0.5ex minus0.5ex}

% margins
\usepackage[tmargin=1in,bmargin=1in,lmargin=1.25in,rmargin=1.25in]{geometry}

% automatically use correct quotes for the current language: \enquote{text in quotes}
\usepackage{csquotes}
\MakeAutoQuote{»}{«}  % input »text in quotes« is equivalent to \enquote{text in quotes}
\let\q=\enquote % use \q{text in quotes}, which is shorter than \enquote{text in quotes}

% erweiterte Einstellungen der Bildunterschriften -> 8 Pt
\usepackage[small]{caption}
\captionsetup{belowskip=12pt,aboveskip=4pt}
\usepackage{subcaption} % group multiple figures together

\usepackage{placeins}
\DeclareOldFontCommand{\sc}{\normalfont\scshape}{\@nomath\sc}
\usepackage[final, nopatch=footnote]{microtype}

% page numbers
\usepackage{tocbasic}
\DeclareTOCStyleEntry[pagenumberwidth=2em]{tocline}{chapter}
\DeclareTOCStyleEntry[pagenumberwidth=2em]{tocline}{section}
\DeclareTOCStyleEntry[pagenumberwidth=2em]{tocline}{subsection}
\DeclareTOCStyleEntry[pagenumberwidth=2em]{tocline}{figure}
\DeclareTOCStyleEntry[pagenumberwidth=2em]{tocline}{table}

%---my Packages-------------------------------------------------------------
\usepackage{tikz}
\usepackage{booktabs}
\usepackage[shortlabels]{enumitem}              % um bessere Aufzählungen zu erstellen
\usepackage{tikz-imagelabels}                   % um Bilder einfach zu beschriften
\usepackage{siunitx}
\sisetup{reset-math-version = false}            % um \boldmath in siunitx zu verwenden
\sisetup{separate-uncertainty = true}           % Unsicherheiten werden mit einem Leerzeichen getrennt
\sisetup{per-mode = fraction}                   % Bruchstrich oder Hochzahl
% Hier Einheiten definieren
\DeclareSIUnit\seppi{sp}                        % Beispiel für eine neue Einheit
\usepackage[style=ieee]{biblatex}               % um Literaturverzeichnisse zu erstellen
\addbibresource{bibliography/references.bib}
\usepackage{longtable}
\usepackage[defaultlines=4,all]{nowidow}        % Verhindert, dass die letzte Zeile eines Absatzes allein auf einer Seite steht
\usepackage{needspace}
\usepackage{acro}                               % um Abkürzungen zu verwalten
\acsetup{
    pages/display=all,
    %single=true,
    single-style=short,
    make-links=true
}

% Grammatikalische Fälle für Abkürzungen (Genitiv, Dativ, Akkusativ).
% Standardmäßig wird im Genitiv ein "s" angehängt (z. B. "des \acs*{pcb}s"),
% Dativ und Akkusativ bleiben unverändert - das deckt die Standardfälle ab,
% ohne dass pro Abkürzung etwas deklariert werden muss. Passt der Standard
% für eine bestimmte Abkürzung nicht, kann er wie bei "plural" gezielt
% überschrieben werden, z.B.:
%   \DeclareAcronym{key}{..., genitive = {}, }                 % kein Suffix
%   \DeclareAcronym{key}{..., short-dative-form = Sonderform, } % eigene Form
\DeclareAcroEnding{genitive}{s}{s}
\DeclareAcroEnding{dative}{}{}
\DeclareAcroEnding{accusative}{}{}

% Befehle analog zu \ac/\acs/\Ac/\Acs, jeweils mit angehängtem Fall:
\NewAcroCommand \acgen   {m}{ \acrogenitive    \UseAcroTemplate {first} {#1} }
\NewAcroCommand \Acgen   {m}{ \acrogenitive    \acroupper \UseAcroTemplate {first} {#1} }
\NewAcroCommand \acsgen  {m}{ \acrogenitive    \UseAcroTemplate {short} {#1} }
\NewAcroCommand \Acsgen  {m}{ \acrogenitive    \acroupper \UseAcroTemplate {short} {#1} }

\NewAcroCommand \acdat   {m}{ \acrodative      \UseAcroTemplate {first} {#1} }
\NewAcroCommand \Acdat   {m}{ \acrodative      \acroupper \UseAcroTemplate {first} {#1} }
\NewAcroCommand \acsdat  {m}{ \acrodative      \UseAcroTemplate {short} {#1} }
\NewAcroCommand \Acsdat  {m}{ \acrodative      \acroupper \UseAcroTemplate {short} {#1} }

\NewAcroCommand \acakk  {m}{ \acroaccusative  \UseAcroTemplate {first} {#1} }
\NewAcroCommand \Acakk  {m}{ \acroaccusative  \acroupper \UseAcroTemplate {first} {#1} }
\NewAcroCommand \acsakk {m}{ \acroaccusative  \UseAcroTemplate {short} {#1} }
\NewAcroCommand \Acsakk {m}{ \acroaccusative  \acroupper \UseAcroTemplate {short} {#1} }

\NewAcroTemplate[list]{symtab}{%
    \AcroNeedPackage{array}%
    \acronymsmapF{%
        \AcroAddRow{%
            \acrowrite{short}&% <-- Das % verhindert den Zeilensprung
            \acrowrite{extra}&%
            \acrowrite{list}\tabularnewline%
        }%
    }%
    {\AcroRerun}%
    \acroheading
    \acropreamble
    \noindent
    \begin{tabular}{m{0.15\linewidth} m{0.15\linewidth} m{0.6\linewidth}}
        \hline
        Symbol & Einheit & Beschreibung \\ \hline
        \AcronymTable
        \hline
    \end{tabular}
}
%---Acronyms-------------------------------------------------------------
\DeclareAcronym{mci}{
	short = 		MCI,
	long =			Management Center Innsbruck,
	tag = 			abbrv,
	first-style =	short}
\DeclareAcronym{ects}{
	short = 		ECTS,
	long = 			European Credit Transfer System,
	tag = 			abbrv,
	first-style =	short}
\DeclareAcronym{WS}{
	short = 		WS,
	long = 			winter semester,
	tag = 			abbrv,
	first-style = 	short}
\DeclareAcronym{SS}{
	short = 		SS,
	long = 			summer semester,
	tag = 			abbrv,
	first-style = 	short}
\DeclareAcronym{etc}{
	short = 		etc.,
	long = 			and so forth,
	tag = 			abbrv,
	first-style = 	short,
	foreign = 		et cetera}
\DeclareAcronym{e.g.}{
	short = 		e.g.,
	long = 			for example,
	tag = 			abbrv,
	first-style = 	short,
	foreign = 		exempli gratia}
\DeclareAcronym{i.e.}{
	short = 		i.e.,
	long = 			in other words,
	tag = 			abbrv,
	first-style = 	short,
	foreign = 		id est}
\DeclareAcronym{vs}{
	short = 		vs.,
	alt = 			vs,
	long = 			versus,
	tag = 			abbrv,
	first-style = 	short}
\DeclareAcronym{approx.}{
	short = 		approx.,
	long = 			approximately,
	tag = 			abbrv,
	first-style = 	short}
\DeclareAcronym{mosfet}{
	short = 		MOSFET,
	long = 			metal-oxide-semulator field-effect transistor,
	tag = 			abbrv,
	first-style = 	short}
\DeclareAcronym{pwm}{
	short = 		PWM,
	long = 			pulse width modulation,
	tag = 			abbrv}
\DeclareAcronym{ac}{
	short = 		AC,
	long = 			alternating current,
	tag = 			abbrv,
	first-style = 	short}
\DeclareAcronym{dc}{
	short = 		DC,
	long = 			direct current,
	tag = 			abbrv,
	first-style = 	short}
\DeclareAcronym{cw}{
	short = 		CW,
	long = 			clockwise,
	tag = 			abbrv}
\DeclareAcronym{ccw}{
	short = 		CCW,
	long = 			counter \acl*{cw},
	tag = 			abbrv}
\DeclareAcronym{ecg}{
	short = 		ECG,
	long = 			electrocardiogram,
	tag = 			abbrv}
\DeclareAcronym{rgb}{
	short = 		RGB,
	long = 			red green blue,
	tag = 			abbrv,
	first-style = 	short}
\DeclareAcronym{led}{
	short = 		LED,
	long = 			light emitting diode,
	tag = 			abbrv,
	first-style = 	short}
\DeclareAcronym{microc}{
	short = 		µC,
	alt = 			uC,
	long = 			microcontroller,
	tag = 			abbrv}
\DeclareAcronym{esp}{
	short = 		ESP,
	long = 			Espressif Systems microcontroller,
	tag = 			abbrv,
	first-style = 	short}
\DeclareAcronym{wifi}{
	short = 		WiFi,
	long = 			Wireless Fidelity,
	tag = 			abbrv,
	first-style = 	short}
\DeclareAcronym{ble}{
	short = 		BLE,
	long = 			Bluetooth Low Energy,
	tag = 			abbrv}
\DeclareAcronym{id}{
	short = 		ID,
	long = 			identification,
	tag = 			abbrv,
	first-style = 	short}
\DeclareAcronym{gpio}{
	short = 		GPIO,
	long = 			general-purpose input/output,
	tag = 			abbrv,
	first-style = 	short}
\DeclareAcronym{cpu}{
	short = 		CPU,
	long = 			central processing unit,
	tag = 			abbrv,
	first-style = 	short}
\DeclareAcronym{rc}{
	short = 		RC,
	long = 			resistor-capacitor,
	tag = 			abbrv,
	first-style = 	short}
\DeclareAcronym{oled}{
	short = 		OLED,
	long = 			organic \acl*{led},
	tag = 			abbrv,
	first-style = 	short}
\DeclareAcronym{3d}{
	short = 		3D,
	long = 			three-dimensional,
	tag = 			abbrv,
	first-style = 	short}
\DeclareAcronym{i2c}{
	short = 		I²C,
	long = 			inter-integrated circuit,
	tag = 			abbrv,
	first-style = 	short}
\DeclareAcronym{sda}{
	short = 		SDA,
	long = 			serial data line,
	tag = 			abbrv,
	first-style = 	short}
\DeclareAcronym{scl}{
	short = 		SCL,
	long = 			serial clock line,
	tag = 			abbrv,
	first-style = 	short}
\DeclareAcronym{spi}{
	short = 		SPI,
	long = 			serial peripheral interface,
	tag = 			abbrv,
	first-style = 	short}
\DeclareAcronym{rf}{
	short = 		RF,
	long = 			radio frequency,
	tag = 			abbrv}
\DeclareAcronym{pcb}{
	short = 		PCB,
	long = 			printed circuit board,
	tag = 			abbrv}
\DeclareAcronym{usb}{
	short = 		USB,
	long = 			universal serial bus,
	tag = 			abbrv,
	first-style = 	short}
\DeclareAcronym{usb-c}{
	short = 		USB-C,
	long = 			\acl*{usb} type C,
	tag = 			abbrv,
	first-style = 	short}
\DeclareAcronym{bom}{
	short = 		BOM,
	long = 			bill of materials,
	tag = 			abbrv}
\DeclareAcronym{imu}{
	short = 		IMU,
	long = 			inertial measurement unit,
	tag = 			abbrv,
	first-style = 	short}
\DeclareAcronym{lipo}{
	short = 		LiPo,
	long = 			lithium polymer,
	tag = 			abbrv,
	first-style = 	short}
\DeclareAcronym{cad}{
	short = 		CAD,
	long = 			computer-aided design,
	tag = 			abbrv,
	first-style = 	short}
\DeclareAcronym{mems}{
	short = 		MEMS,
	long = 			microelectromechanical system,
	tag = 			abbrv,
	short-plural=}
\DeclareAcronym{fr}{
	short = 		FR,
	long = 			functional requirement,
	tag = 			abbrv,
	short-plural=}
\DeclareAcronym{rationale}{
	short = 		R,
	long = 			rationale,
	tag = 			abbrv}
\DeclareAcronym{tc}{
	short = 		TC,
	long = 			test case,
	tag = 			abbrv,
	short-plural=}
\DeclareAcronym{tpu}{
	short = 		TPU,
	long = 			thermoplastic polyurethane,
	tag = 			abbrv,
	first-style = 	short}
\DeclareAcronym{pla}{
	short = 		PLA,
	long = 			polylactic acid,
	tag = 			abbrv,
	first-style = 	short}
\DeclareAcronym{uuid}{
	short = 		UUID,
	long = 			universally unique identifier,
	tag = 			abbrv}
\DeclareAcronym{qol}{
	short = 		QoL,
	long = 			Quality of Life,
	tag = 			abbrv}
\DeclareAcronym{gui}{
	short = 		GUI,
	long = 			Graphical User Interface,
	tag = 			abbrv,
	first-style = 	short}
\DeclareAcronym{fsm}{
	short = 		FSM,
	long = 			finite state machine,
	tag = 			abbrv}



% ---Symbols-------------------------------------------------------------

\DeclareAcronym{acceleration}{
	short = 		$a$,
	long = 			Beschleunigung,
	tag = 			symbol,
	extra = 		$\unit{\frac{m}{s^2}}$
}
\DeclareAcronym{velocity}{
	short = 		$v$,
	long = 			Geschwindigkeit,
	tag = 			symbol,
	extra = 		$\unit{\frac{m}{s}}$
}
\DeclareAcronym{force}{
	short = 		$F$,
	long = 			Kraft,
	tag = 			symbol,
	extra = 		$\unit{N}$
}